% ==========================================================
% 多日多课时排课
% 难度：★★★ 难题
% ==========================================================

\newcommand{\qProbPermMultiDayStem}{%
某年级武术课有太极拳、形意拳、长拳、兵器四门。从周一至周五每天下午安排两门课。每周：太极拳上课三次；形意拳上课三次；长拳上课两次；兵器上课两次；同一门课同一天只上一节。求排课方式数。
}

\newcommand{\qProbPermMultiDayAnswer}{%
$9920$
}

\newcommand{\qProbPermMultiDaySolution}{%
每天下午有两个有先后顺序的课时，五天共十个课时。但同一门课一天至多出现一次。先固定出现次数较多的课程太极拳。

从五天中选三天安排太极拳：$\binom{5}{3}$ 种。在每天，太极拳可位于两个课时之一，有 $2^3$ 种。故太极拳共 $\binom{5}{3}2^3=80$ 种。

太极拳与形意拳都出现三天，至少有一天重合，重合天数只能是 $1, 2, 3$。固定一种太极拳的安排，按重合天数分类：

\textbf{情形一：重合三天}
形意拳必须安排在太极拳所在三天的另一个课时。剩余两个完整日期安排长拳和兵器，每天顺序两种，共 $2^2=4$ 种。

\textbf{情形二：重合两天}
选两天重合：$\binom{3}{2}$ 种。形意拳第三次安排在无太极拳的日期中：$2\times 2=4$ 种位置。此时剩一个完整日期和两个单课时。完整日期填长兵有 $2$ 种，单课时填长兵有 $2$ 种。填法共 $4$ 种。该情形共 $\binom{3}{2}\times 4\times 4 = 48$ 种。

\textbf{情形三：重合一天}
选一天重合：$\binom{3}{1}$ 种。形意拳其余两次在无太极拳的日期：$2^2=4$ 种。剩余四个单课时安排长拳和兵器：$\binom{4}{2}=6$ 种。该情形共 $\binom{3}{1}\times 4\times 6 = 72$ 种。

固定一种太极拳安排后，其余共 $4+48+72=124$ 种。总排课方式数为 $80\times 124 = 9920$ 种。
}

\newcommand{\qProbPermMultiDayTeachingNotes}{%
\textbf{【避坑与边界说明】}
\begin{itemize}
    \item 每天下午的两个课时有先后区别。
    \item 同一门课程同一天不能出现两次，所以不能任意选择十个课时。
    \item 两门各出现三天的课程，在五天内至少重合一天，不能遗漏“重合一天”的必然性。
    \item 分类必须以重合天数为依据，三种情况互斥且完整。
    \item 在重合两天的情形中，填入长拳和兵器时既要考虑完整日期内部的顺序，也要考虑两个单独课时的归属。
\end{itemize}
}
